% =====================================================================
%  BAB 4 — PERANCANGAN SISTEM
% =====================================================================

\chapter{Perancangan Sistem}

\section{Gambaran Umum Sistem}
\label{sec:4.1}

[Berikan gambaran arsitektur sistem secara keseluruhan. Jelaskan komponen-komponen utama,
bagaimana komponen-komponen tersebut berinteraksi, dan bagaimana sistem secara keseluruhan
menjawab rumusan masalah yang ditetapkan.]

\begin{figure}[H]
  \centering
  \gambarbesar{%
    \begin{tikzpicture}[node distance=1.5cm and 2cm]
      % ── Ganti dengan diagram arsitektur Anda ─────────────────────
      \node[procplain, fill=cProsesBg, draw=cProses, text width=3cm]
        (komp1) {Komponen 1};
      \node[procplain, fill=cDataBg, draw=cData, text width=3cm,
            right=of komp1]
        (komp2) {Komponen 2};
      \node[procplain, fill=cOkBg, draw=cOk, text width=3cm,
            right=of komp2]
        (komp3) {Komponen 3};

      \draw[arr] (komp1) -- node[above,font=\scriptsize]{Data} (komp2);
      \draw[arr] (komp2) -- node[above,font=\scriptsize]{Hasil} (komp3);
    \end{tikzpicture}
  }
  \caption{Arsitektur sistem secara keseluruhan}
  \label{fig:4.arsitektur}
\end{figure}

[Deskripsi singkat setiap komponen: peran, input, output.]

\section{Perancangan [Komponen 1]}
\label{sec:4.2}

[Deskripsikan desain komponen pertama secara rinci. Jelaskan algoritma, struktur data,
antarmuka (interface), dan keputusan desain yang diambil beserta alasannya.]

\subsection{[Sub-komponen 1.1]}
\label{sec:4.2.1}

[Detail sub-komponen.]

\subsection{[Sub-komponen 1.2]}
\label{sec:4.2.2}

[Detail sub-komponen.]

\section{Perancangan [Komponen 2]}
\label{sec:4.3}

[Deskripsikan desain komponen kedua.]

\section{Perancangan [Komponen 3]}
\label{sec:4.4}

[Deskripsikan desain komponen ketiga.]

\section{Rancangan Antarmuka Data}
\label{sec:4.5}

[Jelaskan format data yang mengalir antar komponen, skema database (jika ada), format
file input/output, dan protokol komunikasi antar-komponen.]

% Contoh format output JSON
% \begin{lstlisting}[language=,caption={Format data antar-komponen},label={lst:4.format}]
% {
%   "id": "string",
%   "timestamp": "ISO8601",
%   "fitur_a": float,
%   "fitur_b": float,
%   "label": "string"
% }
% \end{lstlisting}

\section{Rancangan Pengujian}
\label{sec:4.6}

[Deskripsikan skenario pengujian yang akan dilakukan pada Bab~6. Apa yang akan diuji,
bagaimana cara mengukurnya, dan kriteria keberhasilan yang ditetapkan.]

\begin{table}[H]
  \centering
  \caption{Skenario pengujian}
  \label{tab:4.skenario}
  \begin{tabular}{clll}
    \toprule
    \textbf{No.} & \textbf{Skenario} & \textbf{Variabel} & \textbf{Metrik} \\
    \midrule
    S1 & [Nama skenario 1] & [Variabel yang diubah] & [Metrik yang diukur] \\
    S2 & [Nama skenario 2] & [Variabel yang diubah] & [Metrik yang diukur] \\
    S3 & [Nama skenario 3] & [Variabel yang diubah] & [Metrik yang diukur] \\
    \bottomrule
  \end{tabular}
\end{table}
